% Compile with: pdflatex supplement_27_2.tex
\documentclass[11pt]{article}
\usepackage[margin=1in]{geometry}
\usepackage{amsmath, amssymb, amsthm}
\usepackage{mathtools}
\usepackage{xcolor}
\usepackage{mdframed}
\usepackage{enumitem}
\usepackage{hyperref}
\usepackage{pagecolor}
\usepackage{microtype}

% ── Colour palette ───────────────────────────────────────────
\definecolor{cream}{HTML}{FFFDF5}
\definecolor{defblue}{HTML}{1A4A7A}
\definecolor{thmgreen}{HTML}{1A5C3A}
\definecolor{warnAmber}{HTML}{7A4A00}
\definecolor{dangerRed}{HTML}{7A1A1A}
\definecolor{boxbluebg}{HTML}{EBF3FB}
\definecolor{boxgreenbg}{HTML}{EBF7EF}
\definecolor{boxamberbg}{HTML}{FDF5E6}
\definecolor{boxredbg}{HTML}{FBEBEB}
\pagecolor{cream}

% ── Theorem styles ───────────────────────────────────────────
\newtheoremstyle{defstyle}  {8pt}{8pt}{\normalfont}{0pt}{\bfseries\color{defblue}} {.}{0.5em}{}
\newtheoremstyle{thmstyle}  {8pt}{8pt}{\normalfont}{0pt}{\bfseries\color{thmgreen}}{.}{0.5em}{}
\newtheoremstyle{notstyle}  {8pt}{8pt}{\normalfont}{0pt}{\bfseries\color{warnAmber}}{.}{0.5em}{}

\theoremstyle{defstyle}
\newtheorem{definition}{Definition}[section]

\theoremstyle{thmstyle}
\newtheorem{theorem}{Theorem}[section]
\newtheorem{corollary}[theorem]{Corollary}

\theoremstyle{notstyle}
\newtheorem*{remark}{Remark}

% ── Framed boxes ─────────────────────────────────────────────
\newmdenv[
  backgroundcolor=boxbluebg, linecolor=defblue, linewidth=2pt,
  topline=false, rightline=false, bottomline=false,
  innerleftmargin=10pt, innerrightmargin=10pt,
  innertopmargin=6pt, innerbottommargin=6pt,
  skipabove=8pt, skipbelow=8pt
]{defbox}

\newmdenv[
  backgroundcolor=boxgreenbg, linecolor=thmgreen, linewidth=2pt,
  topline=false, rightline=false, bottomline=false,
  innerleftmargin=10pt, innerrightmargin=10pt,
  innertopmargin=6pt, innerbottommargin=6pt,
  skipabove=8pt, skipbelow=8pt
]{thmbox}

\newmdenv[
  backgroundcolor=boxamberbg, linecolor=warnAmber, linewidth=2pt,
  topline=false, rightline=false, bottomline=false,
  innerleftmargin=10pt, innerrightmargin=10pt,
  innertopmargin=6pt, innerbottommargin=6pt,
  skipabove=8pt, skipbelow=8pt
]{notebox}

% ── Macros ───────────────────────────────────────────────────
\newcommand{\N}{\mathbb{N}}
\newcommand{\id}{\mathrm{id}}

\begin{document}

% ── Title ────────────────────────────────────────────────────
\begin{center}
  {\Large\bfseries Supplement 27.2 --- $|X \setminus \{a\}| = |X| - 1$}\\[4pt]
  {\normalsize CS 173: Discrete Structures \quad$\cdot$\quad Functions \& Cardinality}\\[2pt]
  {\small\color{gray} University of Illinois at Urbana-Champaign}
\end{center}
\vspace{1em}\hrule\vspace{1.5em}

%=================================================================
\section{Context}
%=================================================================

\begin{notebox}
\begin{remark}[Theorem (!)]
For any finite sets $X$ and $Y$, if $Y \subseteq X$, then
\[
  |X \setminus Y| = |X| - |Y|.
\]
\end{remark}
\end{notebox}

The theorem below is the special case $Y = \{a\}$, which is proved first and used
to establish the general result.

\begin{thmbox}
\begin{theorem}[Theorem (?)]
For any finite set $X$ and any $a$, if $a \in X$, then
\[
  |X \setminus \{a\}| = |X| - 1.
\]
\end{theorem}
\end{thmbox}

%=================================================================
\section{Proof}
%=================================================================

\begin{thmbox}
\begin{proof}
Let $X$ be a finite set, let $a$ be a set, and assume $a \in X$.

Since $X$ is finite, there exists $n \in \N$ such that $|X| = |n|$.
(Recall $n = \{0, 1, 2, \ldots, n-1\}$.)
This means there exists a bijection $f\colon X \to n$; in particular, $f$ is injective
and surjective.

\medskip
\noindent\textbf{Define} $g\colon X \setminus \{a\} \to n-1$ by: for any
$b \in X \setminus \{a\}$,
\[
  g(b) \;:=\;
  \begin{cases}
    f(b)     & \text{if } f(b) < f(a), \\
    f(b) - 1 & \text{if } f(b) > f(a).
  \end{cases}
\]

\begin{notebox}
\begin{remark}
Since $f$ is a bijection and $b \neq a$, we have $f(b) \neq f(a)$, so
$f(b) < f(a)$ or $f(b) > f(a)$ — the two cases are exhaustive and mutually exclusive,
so $g$ is well-defined.
\end{remark}
\end{notebox}

%--- Injectivity --------------------------------------------------
\noindent\textbf{Part 1: Injectivity.}
We prove $(\forall\, b_1, b_2 \in X \setminus \{a\})\bigl(g(b_1) = g(b_2)
\Rightarrow b_1 = b_2\bigr)$.

Let $b_1, b_2 \in X \setminus \{a\}$ and assume $g(b_1) = g(b_2)$.

\medskip
\noindent\underline{Case 1: $f(b_1) < f(a)$.}

By definition of $g$, $g(b_1) = f(b_1)$, so $f(b_1) = g(b_1) = g(b_2)$.

Towards a contradiction, suppose $f(b_2) > f(a)$.
Then by definition, $g(b_2) = f(b_2) - 1 \geq f(a)$ since $f(b_2) > f(a)$.
So $f(b_1) = g(b_1) = g(b_2) \geq f(a)$, contradicting $f(b_1) < f(a)$.
$\Rightarrow\!\Leftarrow$

Therefore $f(b_2) < f(a)$, so by definition $g(b_2) = f(b_2)$.
Thus $f(b_1) = f(b_2)$, so $b_1 = b_2$ because $f$ is injective.

\medskip
\noindent\underline{Case 2: $f(b_1) > f(a)$.}

By definition of $g$, $g(b_1) = f(b_1) - 1$, so $f(b_1) - 1 = g(b_2)$.

Towards a contradiction, suppose $f(b_2) < f(a)$.
Then by definition, $g(b_2) = f(b_2)$, so $f(b_1) - 1 = f(b_2)$.
Since $f(b_1) > f(a)$ we have $f(b_1) \geq f(a) + 1$, hence
$f(b_2) = f(b_1) - 1 \geq f(a)$, contradicting $f(b_2) < f(a)$.
$\Rightarrow\!\Leftarrow$

Therefore $f(b_2) > f(a)$, so by definition $g(b_2) = f(b_2) - 1$.
Thus $f(b_1) - 1 = f(b_2) - 1$, implying $f(b_1) = f(b_2)$,
so $\boxed{b_1 = b_2}$ because $f$ is injective.

\medskip
In both cases $g(b_1) = g(b_2) \Rightarrow b_1 = b_2$, so $g$ is injective.

%--- Surjectivity -------------------------------------------------
\medskip
\noindent\textbf{Part 2: Surjectivity.}
We prove $(\forall\, y \in n{-}1)(\exists\, b \in X \setminus \{a\})\bigl(g(b) = y\bigr)$.

Let $y \in n-1$.  We know $n-1 < n$, so $n-1 \in n$, so $n-1 \subseteq n$,
which gives $y \in n$.

\medskip
\noindent\underline{Case 1: $y < f(a)$.}

Since $f$ is surjective and $y \in n$, there exists $b \in X$ such that $f(b) = y$.

We know $b \neq a$ because if $b = a$, then $y = f(b) = f(a) < f(a)$, a
contradiction. $\Rightarrow\!\Leftarrow$

So $b \in X \setminus \{a\}$ by definition.
Since $f(b) = y < f(a)$, by definition of $g$:
\[
  g(b) = f(b) = y.
\]

\medskip
\noindent\underline{Case 2: $y \geq f(a)$.}

Notice $y + 1 > y \geq f(a)$, so $y + 1 > f(a)$.
Since $y \in n-1$, we know $y < n-1$, so $y+1 < n$, hence $y+1 \in n$.

So, since $f$ is surjective, there exists $b \in X$ such that $f(b) = y+1$.

We know $b \neq a$ because if $b = a$, then $f(a) = f(b) = y+1 > f(a)$, a
contradiction. $\Rightarrow\!\Leftarrow$

So $b \in X \setminus \{a\}$.
Since $f(b) = y+1 > f(a)$, by definition of $g$:
\[
  g(b) = f(b) - 1 = (y+1) - 1 = y.
\]

\medskip
In both cases we found $b \in X \setminus \{a\}$ with $g(b) = y$, so $g$ is
surjective.

\medskip
So, since $g$ is injective and surjective, $g$ is a bijection, and therefore
$|X \setminus \{a\}| = n - 1$ by definition. $\qed$
\end{proof}
\end{thmbox}

\end{document}